\documentclass[11pt, a4paper]{article}

% bfw-style.tex — gemeinsame Style-Definitionen für alle Handouts/Aufgabenhefte
% Voraussetzung: Kompilation mit xelatex (wegen fontspec)

\usepackage[ngerman]{babel}
\usepackage{geometry}
\geometry{a4paper, top=2.2cm, bottom=2.2cm, left=2.3cm, right=2.3cm}

\usepackage{fontspec}
% Fallback-Kette: Source Sans Pro -> Calibri -> Segoe UI -> Latin Modern Sans
% (Latin Modern ist immer mit MiKTeX vorhanden)
\IfFontExistsTF{Source Sans Pro}{%
  \setmainfont{Source Sans Pro}[Ligatures=TeX]%
}{\IfFontExistsTF{Calibri}{%
  \setmainfont{Calibri}[Ligatures=TeX]%
}{\IfFontExistsTF{Segoe UI}{%
  \setmainfont{Segoe UI}[Ligatures=TeX]%
}{%
  \setmainfont{Latin Modern Sans}[Ligatures=TeX]%
}}}
\IfFontExistsTF{Source Code Pro}{%
  \setmonofont{Source Code Pro}[Scale=0.92]%
}{\IfFontExistsTF{Consolas}{%
  \setmonofont{Consolas}[Scale=0.92]%
}{%
  \setmonofont{Latin Modern Mono}[Scale=0.92]%
}}

\usepackage{xcolor}
% Farbpalette: hell, freundlich, modern
\definecolor{bfwPrimary}{HTML}{0EA5A4}    % Petrol/Türkis - Primär
\definecolor{bfwPrimaryDark}{HTML}{0F766E} % dunkler Petrol für Headings
\definecolor{bfwAccent}{HTML}{F59E0B}     % Amber - Akzent
\definecolor{bfwSuccess}{HTML}{10B981}    % Grün - Erfolg / Lösung
\definecolor{bfwWarn}{HTML}{F97316}       % Orange - Achtung
\definecolor{bfwInk}{HTML}{0F172A}        % fast Schwarz
\definecolor{bfwInkSoft}{HTML}{475569}    % gedämpftes Grau
\definecolor{bfwBgSoft}{HTML}{F8FAFC}     % off-white
\definecolor{bfwBgTeal}{HTML}{ECFEFF}     % helles Türkis
\definecolor{bfwBgAmber}{HTML}{FEF3C7}    % helles Gelb
\definecolor{bfwBgRose}{HTML}{FFE4E6}     % helles Rosé für Eselsbrücken

\usepackage[colorlinks=true, linkcolor=bfwPrimaryDark, urlcolor=bfwPrimaryDark]{hyperref}
\usepackage{enumitem}
\setlist[itemize]{leftmargin=*, itemsep=2pt, topsep=2pt}
\setlist[enumerate]{leftmargin=*, itemsep=2pt, topsep=2pt}

\usepackage[most]{tcolorbox}
\usepackage{booktabs}
\usepackage{tikz}
\usetikzlibrary{decorations.pathmorphing, positioning, shapes.geometric, arrows.meta}

% Merkkästchen — wichtige Definition / Kernpunkt
\newtcolorbox{merksatz}[1][]{%
  enhanced, breakable,
  colback=bfwBgTeal,
  colframe=bfwPrimary,
  boxrule=0.6pt,
  arc=4pt,
  left=10pt, right=10pt, top=8pt, bottom=8pt,
  fonttitle=\bfseries\color{bfwPrimaryDark},
  title={\faLightbulb\ Merksatz},
  #1
}

% Eselsbrücke — Mnemoniks
\newtcolorbox{eselsbruecke}[1][]{%
  enhanced, breakable,
  colback=bfwBgRose,
  colframe=bfwAccent,
  boxrule=0.6pt,
  arc=4pt,
  left=10pt, right=10pt, top=8pt, bottom=8pt,
  fonttitle=\bfseries\color{bfwWarn},
  title={Eselsbrücke},
  #1
}

% Beispiel-Box
\newtcolorbox{beispiel}[1][]{%
  enhanced, breakable,
  colback=bfwBgSoft,
  colframe=bfwInkSoft,
  boxrule=0.4pt,
  arc=3pt,
  left=10pt, right=10pt, top=8pt, bottom=8pt,
  fonttitle=\bfseries\color{bfwInk},
  title={Beispiel},
  #1
}

% Lösungs-Box (für Aufgabenhefte)
\newtcolorbox{loesung}[1][]{%
  enhanced, breakable,
  colback=bfwBgAmber!40,
  colframe=bfwSuccess,
  boxrule=0.6pt,
  arc=3pt,
  left=10pt, right=10pt, top=8pt, bottom=8pt,
  fonttitle=\bfseries\color{bfwSuccess!70!black},
  title={Lösung},
  #1
}

% Glossar-Eintrag
\newcommand{\glossareintrag}[2]{%
  \par\noindent\textbf{\color{bfwPrimaryDark}#1} \enspace #2\par\smallskip
}

% Section-Stil — etwas farbiger als Standard
\usepackage{titlesec}
\titleformat{\section}
  {\Large\bfseries\color{bfwPrimaryDark}}
  {\thesection}{0.6em}{}
\titleformat{\subsection}
  {\large\bfseries\color{bfwPrimary}}
  {\thesubsection}{0.5em}{}
\titleformat{\subsubsection}
  {\normalsize\bfseries\color{bfwInk}}
  {\thesubsubsection}{0.4em}{}

% TikZ-Stil "skizzenhaft" — etwas weicher als Rough.js, aber stabil
\tikzset{
  roughbox/.style = {
    draw=bfwPrimaryDark, thick, fill=bfwBgTeal,
    rounded corners=4pt, inner sep=6pt
  },
  roughaccent/.style = {
    draw=bfwAccent, thick, fill=bfwBgAmber,
    rounded corners=4pt, inner sep=6pt
  },
  roughline/.style = {
    draw=bfwInkSoft, thick, -{Stealth[length=2.5mm]}
  }
}

% Bullet-Symbole (bewusst kein FontAwesome — vermeidet Font-Installations-Probleme)
\newcommand{\faLightbulb}{$\bullet$}
\newcommand{\faBookmark}{$\bullet$}

% Header/Footer
\usepackage{fancyhdr}
\pagestyle{fancy}
\fancyhf{}
\renewcommand{\headrulewidth}{0pt}
\fancyfoot[L]{\small\color{bfwInkSoft}\thepage}
\fancyfoot[R]{\small\color{bfwInkSoft} BFW M-064 Beschaffung im Büromanagement}


\title{\vspace*{-1.5cm}\Huge\bfseries\color{bfwPrimaryDark}Tag~6: Kaufvertragsstörungen\\[0.4em]
  \large\color{bfwInkSoft}\textnormal{Lieferungsverzug, Annahmeverzug, Schlechtleistung, Käuferrechte}}
\author{\textsf{Beschaffung im Büromanagement}\\
\textsf{\small\copyright\ Can Siebert\,\textperiodcentered\,2026}}
\date{19.05.2026}

\begin{document}
\maketitle
\thispagestyle{empty}

\vspace{1em}
\begin{merksatz}[title={\bfseries Worum geht es heute?}]
Im Idealfall läuft der Kaufvertrag glatt: Ware wird geliefert, angenommen, bezahlt.
In der Praxis gibt es jedoch häufig \emph{Leistungsstörungen} --- der Lieferant ist
zu spät, die Käuferin nimmt nicht ab, die Ware ist mangelhaft. Heute lernen wir die
drei Hauptstörungen, die gesetzlichen Grundlagen (§§\,286, 293\,ff., 437 BGB) und
welche Rechte und Pflichten daraus für beide Seiten entstehen --- mit einem klaren
Blick auf die IHK-Prüfung Teil~1 und Teil~2.
\end{merksatz}

\tableofcontents
\vspace{1em}
\hrule
\vspace{1em}

\section{Überblick: Leistungsstörungen}

Eine \emph{Leistungsstörung} liegt vor, wenn die geschuldete Leistung nicht, nicht
rechtzeitig oder nicht ordnungsgemäß erbracht wird. Im Kaufrecht unterscheiden wir
drei Hauptarten:

\begin{itemize}[itemsep=2pt]
  \item \textbf{Schuldnerverzug} (Lieferungsverzug oder Zahlungsverzug) --- die Leistung kommt zu spät.
  \item \textbf{Gläubigerverzug} (Annahmeverzug) --- die Leistung wird angeboten, aber nicht angenommen.
  \item \textbf{Schlechtleistung} (mangelhafte Lieferung) --- es wird geliefert, aber die Ware ist mangelhaft.
\end{itemize}

\begin{eselsbruecke}
\textbf{,,SAS"} hilft beim Merken: \textbf{S}chuldnerverzug --- \textbf{A}nnahmeverzug --- \textbf{S}chlechtleistung. Die drei großen Familien der Kaufvertragsstörungen.
\end{eselsbruecke}

\section{Lieferungsverzug --- § 286 BGB}

\subsection{Voraussetzungen des Schuldnerverzugs}

Damit ein Schuldner (z.\,B.\ der Lieferant) in Verzug kommt, müssen drei Voraussetzungen
erfüllt sein:

\begin{enumerate}[itemsep=2pt]
  \item \textbf{Fälligkeit der Leistung} --- der Liefertermin ist erreicht.
  \item \textbf{Mahnung} durch den Gläubiger --- nach Eintritt der Fälligkeit.
  \item \textbf{Vertretenmüssen} des Schuldners --- er hat den Verzug zu verschulden (Vorsatz oder Fahrlässigkeit, § 276 BGB). Wird vermutet.
\end{enumerate}

\subsection{Wann ist die Mahnung entbehrlich?}

Nach §\,286 Abs.\,2 BGB braucht es \emph{keine} Mahnung, wenn:

\begin{itemize}[itemsep=2pt]
  \item ein \textbf{kalendermäßig bestimmter} Liefertermin vereinbart ist (,,Lieferung am 15.\,Mai 2026").
  \item ein Termin sich \textbf{kalendermäßig berechnen} lässt (,,Lieferung 14 Tage nach Bestellung").
  \item der Schuldner die Leistung \textbf{ernsthaft und endgültig verweigert}.
  \item besondere Gründe einen sofortigen Verzugseintritt rechtfertigen.
\end{itemize}

\begin{merksatz}[title={Praxisregel}]
Bei Geldschulden gilt zudem: Spätestens 30 Tage nach Fälligkeit und Zugang einer
Rechnung kommt der Schuldner ohne Mahnung in Verzug (§\,286 Abs.\,3 BGB).
\end{merksatz}

\subsection{Rechtsfolgen des Schuldnerverzugs}

Ist der Schuldner im Verzug, hat der Gläubiger folgende Ansprüche:

\begin{itemize}[itemsep=2pt]
  \item \textbf{Erfüllung} bleibt bestehen --- die Lieferung muss nachgeholt werden.
  \item \textbf{Schadensersatz wegen Verzögerung} (§§ 280 I, 286 BGB) --- z.\,B.\ entgangener Gewinn.
  \item \textbf{Verzugszinsen} (§ 288 BGB).
  \item \textbf{Rücktritt} nach Setzen einer angemessenen Nachfrist (§ 323 BGB).
  \item \textbf{Schadensersatz statt der Leistung} (§ 281 BGB) nach Fristablauf.
\end{itemize}

\subsection{Verzugszinsen --- die wichtige Rechnung}

\begin{merksatz}[title={Formel}]
\textbf{Verzugszinsen = Forderung $\times$ Zinssatz $\times$ Tage $\div$ 360}
\end{merksatz}

Der \emph{Verzugszinssatz} unterscheidet zwei Fälle:

\begin{center}
\begin{tabular}{p{6cm} p{8cm}}
\toprule
\textbf{Fall} & \textbf{Zinssatz}\\
\midrule
Geschäfte mit Verbrauchern (B2C) & Basiszinssatz $+$ 5\,\%-Punkte\\
Geschäfte zwischen Unternehmern (B2B, kein Verbraucher beteiligt) & Basiszinssatz $+$ 9\,\%-Punkte\\
\bottomrule
\end{tabular}
\end{center}

Der Basiszinssatz wird halbjährlich von der Deutschen Bundesbank festgelegt
(aktueller Wert z.\,B.\ 3{,}62\,\%, kann sich ändern).

\begin{beispiel}
Eine Rechnung über 8.000~€ ist seit 30~Tagen nicht beglichen --- B2B-Verhältnis,
Basiszins 3{,}62\,\%. Verzugszinssatz: $3{,}62 + 9 = 12{,}62\,\%$. Verzugszinsen
$= 8\,000 \times 0{,}1262 \times \frac{30}{360} = \mathbf{84{,}13~€}$.
\end{beispiel}

\section{Annahmeverzug --- §§ 293 ff.\ BGB}

\subsection{Voraussetzungen}

Der Gläubiger (z.\,B.\ Käufer) gerät in Annahmeverzug, wenn:

\begin{enumerate}[itemsep=2pt]
  \item Der Schuldner die Leistung \textbf{ordnungsgemäß angeboten} hat (Zeit, Ort, Art, § 294 BGB --- tatsächliches Angebot, oder § 295 BGB --- wörtliches Angebot bei Mitwirkung des Gläubigers).
  \item Der Gläubiger die Leistung \textbf{nicht angenommen} hat oder eine erforderliche Mitwirkung unterlässt.
\end{enumerate}

\subsection{Rechtsfolgen für den Schuldner}

\begin{itemize}[itemsep=2pt]
  \item \textbf{Haftungsmilderung} (§ 300 BGB) --- der Schuldner haftet nur noch für Vorsatz und grobe Fahrlässigkeit.
  \item \textbf{Mehraufwendungen} (z.\,B.\ Lager- und Transportkosten) sind zu ersetzen (§ 304 BGB).
  \item \textbf{Gefahrübergang} bei Stückschulden auf den Gläubiger (§ 300 BGB iVm § 446 BGB).
\end{itemize}

\begin{eselsbruecke}
,,Wer nicht annimmt, trägt das Risiko." Sobald der Käufer in Annahmeverzug ist,
muss er auch die zusätzlichen Kosten und das Risiko mittragen.
\end{eselsbruecke}

\section{Schlechtleistung --- § 437 BGB}

\subsection{Was ist ein Sachmangel?}

Ein \emph{Sachmangel} liegt nach § 434 BGB vor, wenn die Sache bei Gefahrübergang
nicht die vereinbarte Beschaffenheit hat, sich nicht für die vertraglich vorausgesetzte
Verwendung eignet oder nicht die übliche Beschaffenheit aufweist, die der Käufer
erwarten darf.

Auch eine \textbf{Falschlieferung} (aliud) und eine \textbf{Mengenabweichung} (Minus)
gelten nach § 434 III BGB als Sachmangel.

\subsection{Die vier Käuferrechte}

\begin{merksatz}[title={§ 437 BGB --- Käuferrechte im Überblick}]
\textbf{1.\ Nacherfüllung} (Vorrang!) \textbar\ \textbf{2.\ Rücktritt} \textbar\ \textbf{3.\ Minderung} \textbar\ \textbf{4.\ Schadensersatz}
\end{merksatz}

\textbf{1.\ Nacherfüllung (§ 439 BGB) --- der Vorrang.} Der Käufer muss zuerst dem
Verkäufer die Möglichkeit zur Nacherfüllung geben. Er hat die Wahl zwischen
\emph{Nachbesserung} (Reparatur) und \emph{Ersatzlieferung} (mangelfreie Sache).

\textbf{2.\ Rücktritt (§§ 437 Nr.\,2, 323 BGB).} Erst nach erfolglosem Ablauf einer
\emph{angemessenen Frist} darf der Käufer zurücktreten. Bei unerheblichen Mängeln ist
der Rücktritt ausgeschlossen (§ 323 V 2 BGB).

\textbf{3.\ Minderung (§ 441 BGB).} Statt zurückzutreten kann der Käufer den Kaufpreis
\emph{mindern}. Die Minderung ist auch bei unerheblichen Mängeln möglich.

\textbf{4.\ Schadensersatz (§§ 437 Nr.\,3, 280, 281 BGB).} Zusätzlich zu Rücktritt
oder Minderung möglich.

\subsection{Vorrang der Nacherfüllung --- warum?}

Der Gesetzgeber will dem Verkäufer die Chance geben, den Mangel zu beheben, bevor
der Käufer schwerwiegendere Rechte ausübt. Erst wenn die Nacherfüllung scheitert,
unmöglich oder unzumutbar ist, sind Rücktritt, Minderung und Schadensersatz statt
der Leistung möglich.

\begin{merksatz}[title={Angemessene Frist}]
Die Rechtsprechung akzeptiert in der Regel \textbf{14 Tage} als angemessen --- bei
einfachen Reparaturen kann auch weniger ausreichen, bei komplizierten Maschinen mehr.
\end{merksatz}

\subsection{Minderungsformel}

\begin{merksatz}[title={Formel nach § 441 Abs. 3 BGB}]
\textbf{Geminderter Preis = Kaufpreis $\times \dfrac{\text{Wert mangelhaft}}{\text{Wert mangelfrei}}$}
\end{merksatz}

\begin{beispiel}
Drucker für 600~€ gekauft. Wert mangelfrei: 600~€. Wert mangelhaft: 420~€.
Geminderter Preis $= 600 \times \frac{420}{600} = \mathbf{420~€}$. Rückforderung: 180~€.
\end{beispiel}

\section{Besonderheiten der Kaufvertragsarten}

\subsection{Handelskauf}

Beim \emph{Handelskauf} (Kauf zwischen Kaufleuten) gelten zusätzlich die §§ 373 ff.\ HGB.
Wichtigste Besonderheit: die \textbf{Rügepflicht} nach § 377 HGB. Der Käufer muss die
Ware unverzüglich untersuchen und Mängel sofort anzeigen --- sonst gilt die Ware als
genehmigt und alle Mängelrechte sind verloren.

\subsection{Verbrauchsgüterkauf}

Verbrauchsgüterkauf (§ 474 BGB) ist der Kauf einer beweglichen Sache durch einen
\emph{Verbraucher} (§ 13 BGB) von einem \emph{Unternehmer} (§ 14 BGB). Die
Käuferrechte sind hier zugunsten des Verbrauchers verstärkt:

\begin{itemize}[itemsep=2pt]
  \item \textbf{Beweislastumkehr} (§ 477 BGB): In den ersten 12~Monaten nach Übergabe wird vermutet, dass ein auftretender Mangel bereits bei Übergabe vorlag.
  \item \textbf{Keine Verkürzung} der Verjährungsfrist auf unter zwei Jahre bei Neuware.
  \item \textbf{Garantieerklärungen} müssen klar und transparent sein (§ 479 BGB).
\end{itemize}

\subsection{Stückkauf vs.\ Gattungskauf}

\begin{center}
\begin{tabular}{p{4cm} p{5cm} p{5cm}}
\toprule
\textbf{Kriterium} & \textbf{Stückkauf} & \textbf{Gattungskauf}\\
\midrule
Bestimmung & Konkrete, einzelne Sache & Gattungsbezeichnung (z.\,B.\ ,,100 Aktenordner blau")\\
Beispiel Bürokontext & Gebrauchter Drucker mit Seriennr. XYZ & 50 neue Tonerkartuschen\\
Untergang & Vertrag wird unmöglich, falls die einzige Sache untergeht & Lieferant muss aus der Gattung neu beschaffen\\
\bottomrule
\end{tabular}
\end{center}

\section{Zusammenfassung}

\begin{enumerate}[itemsep=2pt]
  \item \textbf{Drei Störungen}: Schuldner-, Gläubigerverzug, Schlechtleistung.
  \item \textbf{§ 286 BGB}: Fälligkeit + Mahnung + Vertretenmüssen $=$ Verzug. Mahnung entbehrlich bei kalendermäßigem Termin.
  \item \textbf{Verzugszinsen}: B2B Basiszins $+$ 9\,\%-Punkte; B2C $+$ 5\,\%-Punkte.
  \item \textbf{§§ 293\,ff.\ BGB}: Annahmeverzug $\to$ Haftungsmilderung, Mehraufwendungen, Gefahrübergang.
  \item \textbf{§ 437 BGB}: Nacherfüllung (Vorrang) $\to$ Rücktritt/Minderung $\to$ Schadensersatz.
  \item \textbf{Verbrauchsgüterkauf}: Beweislastumkehr 12 Monate (§ 477 BGB).
  \item \textbf{Stückkauf vs.\ Gattungskauf}: konkret vs.\ aus der Gattung.
\end{enumerate}

\newpage

\section*{Glossar}
\addcontentsline{toc}{section}{Glossar}

\glossareintrag{Leistungsstörung}{Sammelbegriff für alle Fälle, in denen die geschuldete Leistung nicht, nicht rechtzeitig oder nicht ordnungsgemäß erbracht wird.}

\glossareintrag{Schuldnerverzug}{Verzug des Leistungsschuldners (z.\,B.\ Lieferanten) nach § 286 BGB --- Voraussetzungen: Fälligkeit, Mahnung, Vertretenmüssen.}

\glossareintrag{Lieferungsverzug}{Schuldnerverzug bezüglich der Lieferungspflicht --- der Lieferant ist zu spät.}

\glossareintrag{Zahlungsverzug}{Schuldnerverzug bezüglich der Zahlungspflicht --- der Käufer zahlt verspätet (relevant in Tag~8).}

\glossareintrag{Mahnung}{Eindeutige Aufforderung des Gläubigers an den Schuldner, die fällige Leistung zu erbringen --- in der Regel schriftlich.}

\glossareintrag{Kalendermäßige Bestimmung}{Liefertermin nach dem Kalender festgelegt (,,15.05.2026") oder berechenbar (,,14 Tage nach Bestellung") --- Mahnung entbehrlich.}

\glossareintrag{Vertretenmüssen}{Verschulden des Schuldners (Vorsatz oder Fahrlässigkeit, § 276 BGB) --- wird beim Verzug vermutet.}

\glossareintrag{Verzugszinsen}{Gesetzlicher Zinsanspruch des Gläubigers nach § 288 BGB --- B2B Basiszins $+$ 9\,\%-Punkte, B2C $+$ 5\,\%-Punkte.}

\glossareintrag{Basiszinssatz}{Halbjährlich von der Deutschen Bundesbank festgesetzter Referenzzinssatz, Basis für Verzugszinsen.}

\glossareintrag{Annahmeverzug}{Verzug des Gläubigers (§§ 293\,ff.\ BGB), wenn er die ordnungsgemäß angebotene Leistung nicht annimmt.}

\glossareintrag{Haftungsmilderung}{Folge des Annahmeverzugs (§ 300 BGB) --- der Schuldner haftet nur noch für Vorsatz und grobe Fahrlässigkeit.}

\glossareintrag{Sachmangel}{Abweichung der Sache von der vereinbarten oder üblichen Beschaffenheit (§ 434 BGB) --- inklusive Falschlieferung und Mengenabweichung.}

\glossareintrag{Nacherfüllung}{Vorrangiges Käuferrecht (§ 439 BGB): Wahl zwischen Nachbesserung (Reparatur) und Ersatzlieferung.}

\glossareintrag{Rücktritt}{Käuferrecht (§§ 437 Nr.\,2, 323 BGB) nach erfolgloser Fristsetzung --- der Vertrag wird rückabgewickelt.}

\glossareintrag{Minderung}{Käuferrecht (§ 441 BGB) --- Herabsetzung des Kaufpreises im Verhältnis Wert mangelfrei zu Wert mangelhaft.}

\glossareintrag{Schadensersatz statt der Leistung}{Käuferrecht (§ 281 BGB), wenn Nacherfüllung scheitert --- der Käufer kann den Vertrag aufheben und Schadensersatz fordern.}

\glossareintrag{Angemessene Frist}{Zeitraum, den der Käufer dem Verkäufer für die Nacherfüllung setzt --- in der Regel 14~Tage.}

\glossareintrag{Handelskauf}{Kaufvertrag zwischen Kaufleuten --- es gelten zusätzlich §§ 373\,ff.\ HGB, insbesondere § 377 HGB Rügepflicht.}

\glossareintrag{Verbrauchsgüterkauf}{Kauf einer beweglichen Sache durch einen Verbraucher von einem Unternehmer (§ 474 BGB) --- mit Schutzregeln zugunsten des Käufers.}

\glossareintrag{Beweislastumkehr}{§ 477 BGB --- in den ersten 12~Monaten nach Übergabe wird vermutet, dass ein Mangel bereits bei Übergabe vorlag.}

\glossareintrag{Stückkauf}{Kauf einer konkreten, individuell bestimmten Sache (z.\,B.\ ein gebrauchter Drucker mit Seriennummer).}

\glossareintrag{Gattungskauf}{Kauf einer der Gattung nach bestimmten Sache (z.\,B.\ ,,50 Aktenordner blau") --- Lieferant kann aus der Gattung beschaffen.}

\end{document}
